\documentclass[11pt]{article}
\usepackage[letterpaper,margin=1in]{geometry}
\usepackage{booktabs}
\usepackage{amsmath}
\usepackage{enumitem}
\usepackage{hyperref}
\usepackage[T1]{fontenc}
\usepackage{lmodern}
\title{\textbf{Crossed Control in China, 1944--1945}\\
\large Applying the \emph{Battle for Germany / Downfall} mechanism to the defeat of Japan and the coming Chinese Civil War}
\author{}
\date{}
\begin{document}
\maketitle

\begin{center}
\begin{tabular}{lll}
\toprule
Player & Primary faction & Cross-controlled opponent \\
\midrule
CCP player & CCP forces & Japanese forces principally confronting the KMT \\
KMT player & KMT/NRA forces & Japanese forces principally confronting the CCP \\
\bottomrule
\end{tabular}
\end{center}

\section{Design Premise}
China in 1944-45 offers a close strategic analogue to the end of the European war: two nominal allies have a common immediate enemy, but each is already positioning itself for the political and military contest that will follow the enemy's defeat.


The recommendation is not to copy the geographical split of Battle for Germany or Downfall literally. Germany permits a clean east-west division. China instead calls for crossed control based on political opportunity cost: each player controls his own Chinese faction and the Japanese forces whose principal task is to obstruct the rival Chinese faction.


\textbf{You command your Chinese faction against Japan, and you command the Japanese resistance that obstructs your Chinese rival.}


Neither player wants Japan to win the war. Each nevertheless has a strong incentive to make the Japanese forces facing the other Chinese faction fight intelligently and tenaciously.


\section{Victory: The Postwar Position}
Japanese defeat should end the common war but should not by itself determine the winner. Victory should instead depend on each Chinese faction's position when Japan surrenders.


\textbf{P = T + C + A + L - D}


The terms can represent politically valuable territory or population (T), major cities and communications (C), surviving armed strength (A), legitimacy or political organization (L), and disruption, overextension, or dependence (D).


The KMT and CCP need not use identical weights. The KMT should care disproportionately about recognized governmental authority, major cities, communications, and provincial capitals. The CCP should place greater value on rural political control, base areas, population, cadres, and intact military organization.


\textbf{Japan's defeat ends the war, but the condition in which each Chinese faction reaches that moment determines who has won the game.}


\section{Japan Must Have an Operational Agenda}
Cross-control fails if a player can simply use his Japanese formations to sit passively in front of the rival. Japan therefore requires independent strategic constraints. Spring 1944 is an attractive starting point because Operation Ichi-Go supplies a strong opening agenda.


Japanese formations should have mandatory or heavily incentivized objectives: open and maintain the north-south railway corridor; capture or neutralize air bases; hold major cities and communications centers; protect railways and supply lines; and suppress military threats in occupied areas.


The opposing Chinese player chooses how Japan pursues those legitimate objectives. The CCP player controlling Japan against the KMT might choose which Nationalist army receives the main blow, whether to exploit a breakthrough or consolidate, whether destruction of KMT formations is preferable to territorial gain, and where scarce Japanese support should be concentrated.


Every available choice should be militarily plausible for Japan, while naturally allowing the controlling player to prefer the alternative that leaves his Chinese rival in the weakest postwar position.


\section{A Four-Actor Time Economy}
A useful extension from Downfall would be an initiative or time-track system with four operational actors: \textbf{KMT - CCP - Japan/KMT Front - Japan/CCP Front.}


An operation consumes time or initiative; the actor furthest behind on the track becomes eligible to act. This makes operational tempo a resource rather than imposing a conventional alternating turn.


China suggests an additional mechanism: Japanese strategic resources should be shared. The two Japanese commands might draw from common pools of replacements, air support, logistics, transport, reserves, or operational capacity.


If the CCP player spends Japanese resources on another major offensive against the KMT, those resources are unavailable to the Japanese forces that the KMT player is using against the CCP.


Each player is simultaneously trying to defeat Japan with his Chinese faction; preserve and strengthen that faction for the coming civil war; use Japanese forces to damage his rival; and compete over the same Japanese war machine while controlling different portions of it.


\section{Japanese Surrender as an Endgame}
Germany in Downfall is conquered by advancing armies. Japan's surrender in China creates a different problem: the decisive question becomes who occupies the territory, communications, arms stocks, and political vacuum left by surrendering Japanese forces.


Japanese surrender should therefore trigger a short Endgame Phase rather than immediately ending play.


When an external Pacific War track reaches its surrender threshold: Japanese offensive combat ceases; occupied cities, rail junctions, depots, and surrendering formations become strategic prizes; and the CCP and KMT receive perhaps two or three special initiative rounds in which to occupy or establish control over those assets.


Earlier decisions then acquire their proper consequences. A player who has used Japanese forces effectively against his rival may have left that rival too weak to exploit the surrender. Conversely, excessive concentration on weakening the rival may have allowed the opponent's Japanese command to devastate one's own Chinese forces.


\section{Manchuria as the Berlin Equivalent}
Manchuria should be a separate endgame subsystem. Soviet entry is best treated as an exogenous event rather than placing the Soviet armies under either player's direct control.


When Soviet intervention occurs, the Kwantung Army can rapidly collapse, become immobilized, or otherwise cease to function as a conventional Japanese field force. Japanese depots, equipment, communications, and territory then become potential assets in the postwar struggle.


The CCP values access, cadres, bases, population, and military resources in Manchuria; the KMT values recognized governmental control, major cities, railways, and communications.


The campaign can consequently evolve from a war against Japan into a contest over the shape of postwar China without requiring unrestricted KMT-CCP combat during most of the game.


\section{Chinese-on-Chinese Combat and Legitimacy}
Direct KMT-CCP combat should be possible but politically expensive. A United Front, Legitimacy, or similar track can impose costs for openly attacking the other Chinese faction. Such attacks might reduce political standing, foreign support, victory points, or access to resources.


\textbf{Using one's own Chinese troops to destroy the rival is politically expensive. Using Japanese forces to damage that rival is not.}


Cross-control then ceases to be merely a convenient player-assignment device. It becomes a model of the strategic incentive to welcome damage to a future domestic opponent while still requiring the common enemy ultimately to be defeated.


\section{Dynamic Rather Than Permanent Japanese Control}
A literal geographical division of Japanese forces is less appropriate to China than to Germany. Japanese formations should instead be assigned according to the opponent against whom they are principally operating.


\textbf{J(KMT) $\rightarrow$ controlled by the CCP player}\\\textbf{J(CCP) $\rightarrow$ controlled by the KMT player}


A Japanese formation transferred from anti-CCP operations to a major campaign against Nationalist forces could therefore change player control.


Neither player is Japan. Japan remains a third actor whose local operational decision-making is delegated to whichever player has the strongest incentive to make that resistance effective.


Safeguards would be needed to prevent deliberate abuse during a control transfer: restrictions on self-sabotaging movement immediately before reassignment, mandatory Japanese objectives, or a neutral strategic directive governing transfers.


\section{The Central Strategic Question: How Quickly Should Japan Lose?}
The most promising emergent decision is not simply how to defeat Japan, but \textbf{how quickly each player wants Japan to lose.}


A player whose faction currently holds the superior postwar position wants surrender soon. A player whose rival is better positioned may want another operational cycle - not because he wants Japan to win, but because the Japanese forces he controls may still have time to damage that rival before disappearing as an active belligerent.


The desired duration of the common enemy thus becomes endogenous to the game state.


\textbf{The game begins as a struggle to defeat Japan and gradually reveals itself as a struggle to win the Chinese Civil War before that war officially resumes.}


\section{Recommended Prototype Structure}
1. Start in spring 1944, with Operation Ichi-Go supplying the initial Japanese strategic agenda.


2. Use two human players: KMT and CCP.


3. Give each player operational control of Japanese forces primarily opposing the other Chinese faction.


4. Constrain Japanese actions through objectives, command geography, and a shared Japanese resource economy.


5. Use four initiative/time markers rather than a conventional I-go-you-go turn.


6. Score asymmetric KMT and CCP postwar objectives rather than Japanese casualties alone.


7. Penalize overt Chinese-on-Chinese combat through legitimacy or United Front costs.


8. Allow Japanese formations to change player control when their principal operational opponent changes.


9. Trigger Soviet entry and Japanese surrender through an external Pacific War track.


10. Follow surrender with a short occupation/race endgame, with Manchuria receiving special treatment.


The essential design test is whether players repeatedly face plausible choices in which defeating Japan, preserving their own faction, weakening their future rival, and determining the timing of surrender point in different directions. If those tensions arise without arbitrary exceptions, crossed control is doing genuine historical work rather than serving merely as a novelty.


\end{document}
